\documentclass[11pt]{article}
\usepackage[margin=25mm]{geometry}
\usepackage{fontspec}
\setmainfont{Latin Modern Roman}
\usepackage{amsmath,amssymb,mathtools}
\usepackage{unicode-math}
\setmathfont{Latin Modern Math}
\usepackage{xurl,hyperref}
\hypersetup{colorlinks=true,urlcolor=blue}
\setcounter{secnumdepth}{0}
\setlength{\parindent}{0pt}
\setlength{\parskip}{5pt}
\setlength{\emergencystretch}{3em}
\newcommand{\sref}[2]{\hyperlink{source:#1:#2}{[S#2]}}
\newcommand{\cataloguescope}[1]{\paragraph{Recorded target.}#1}
\begin{document}
\section{Research Problems in Function Theory — Problem 6.106}\label{statement}

\subsection{Definitions and mathematical statement}

\paragraph{Shared notation.}$\mathbb N=\{1,2,\ldots\}$, and $\mathbb Z,\mathbb Q,\mathbb R,\mathbb C$ denote the integers, rationals, reals and complex numbers. Any other conventions are specified locally.


Write $\mathbb D=\{z\in\mathbb C:|z|<1\}$ and $\mathbb T=\partial\mathbb D$. Let $J$ be a Jordan curve, that is, the image of the circle under a continuous injective map, and let $D$ be its bounded complementary domain. Let $b:\mathbb T\to J$ be an orientation-preserving homeomorphism with respect to the positive boundary orientations. Its harmonic extension is the Poisson integral
\[u(re^{i\theta})=\frac1{2\pi}\int_0^{2\pi}
\frac{1-r^2}{1-2r\cos(\theta-t)+r^2}\,b(e^{it})\,dt,\qquad0\leq r<1.\]
\paragraph{Statement recorded in the supplied source.}
If $u(\mathbb D)\subseteq D$, must $u:\mathbb D\to D$ be a homeomorphism onto $D$? No smoothness of the Jordan curve or boundary homeomorphism is imposed. \sref{statement}{2}
\subsection{Short English statement}
Does a harmonic extension of an oriented Jordan-boundary homeomorphism become a homeomorphism whenever its entire image stays inside the target domain?
\subsection{Sources}
\begin{description}
\item[\hypertarget{source:statement:1}{S1}] ULAM.AI and the UnsolvedMath Contributors, UnsolvedMath: A Curated Collection of Open Mathematics Problems, catalogue entry. CC BY 4.0. \url{https://huggingface.co/datasets/ulamai/UnsolvedMath}.
\item[\hypertarget{source:statement:2}{S2}] W. K. Hayman and E. F. Lingham, Research Problems in Function Theory (2018), Problem 6.106, complete problem and update. \url{https://arxiv.org/abs/1809.07200}.
\end{description}
\label{end:statement}
\end{document}
